\documentclass[11pt,a4paper]{article}
\usepackage[utf8]{inputenc}
\usepackage{amsmath,amssymb,amsfonts,amsthm}
\usepackage{geometry}
\geometry{margin=1in}
\usepackage{booktabs}
\usepackage{tabularx}
\usepackage{hyperref}
\usepackage{microtype}
\emergencystretch=2em

\hypersetup{colorlinks=true,linkcolor=blue,citecolor=blue,urlcolor=blue}

\newtheorem{proposition}{Proposition}
\newtheorem{corollary}[proposition]{Corollary}
\theoremstyle{definition}
\newtheorem{remark}[proposition]{Remark}

\title{\textbf{Appendix A: Equilibrium and Linear-Stability Analysis of the
Physics v1.0 Dynamical Sector}\\[0.4em]
\large Supplementary Artifact v1.1 --- Additive to Physics v1.0}

\author{
    \textbf{Prince Upadhyay}\\
    \small Independent Researcher\\
    \small Email: \texttt{starxshark@gmail.com}
}

\date{\today}

\begin{document}
\maketitle

\begin{abstract}
Physics v1.0 introduces a three-variable phenomenological dynamical system and
establishes its dimensional closure, but does not analyse it. This appendix
supplies that analysis. We show that the system possesses a unique equilibrium
in closed form; that the steady-state processing rate is fixed by input power
alone and is independent of the entire entropy sector; that the equilibrium
distance from the informational saturation parameter is set exactly by the
irreversible-erasure rate; and that the system's only linear instability is a
Hopf bifurcation whose threshold scales as the inverse of that same erasure
rate, so that the system is unconditionally linearly stable in the
zero-erasure limit. We also record two constraints absent from v1.0: an
admissibility bound on the erasure rate, and a finite-time escape to
unphysical states that requires an explicit domain restriction. All analytic
results are verified numerically to machine precision.

This appendix is \textbf{additive}. It introduces no new postulate, modifies no
equation, and has no bearing on the Gate~0 verdict.
\end{abstract}

\tableofcontents
\newpage

\section{Scope and status}

This document is a supplementary artifact. Under the immutability rules of the
frozen research lineage:

\begin{enumerate}
\item Physics v1.0 remains v1.0. Nothing here edits it.
\item No equation, coefficient, or interpretation of v1.0 is altered. The
system analysed below is exactly equations (9)--(11) of
\texttt{paper\_physics\_v1.0\_release.tex} as frozen.
\item No new postulate is introduced. Every statement is a consequence of the
frozen ansatz.
\item Gate~0 Technical Report v1.0 and its verdict \textbf{UNRESOLVED IN
GENERALITY} are untouched. Gate~0 audits the observer-entropy to
horizon-entropy bridge; the present appendix analyses only the dynamical
sector, which Gate~0 did not examine.
\end{enumerate}

\begin{remark}[$H_{\max}$ carries no geometric content here]
\label{rem:hmax}
Throughout this appendix $H_{\max}$ is treated strictly as the
\emph{informational saturation parameter in bits} declared in v1.0's dynamical
sector. Its identification with a causal-horizon capacity
$A/(4\ell_P^2\ln 2)$ is the structural domain extension whose microphysical
justification v1.0 leaves open and Gate~0 finds unresolved. No result below
depends on that identification, and none of them should be read as supporting
it.
\end{remark}

\section{The system under analysis}

Write $\dot I \equiv \dot I_{\text{model}}$ and $\dot I_e \equiv
\dot I_{\text{erase}}$. The frozen system is
\begin{align}
\frac{d\gamma}{dt} &= a\,\dot I\,\mathcal{E}_{\text{avail}}
  \left(1-\frac{\gamma}{\gamma_{\max}}\right)-b\gamma,
  \label{eq:gamma}\\[4pt]
\frac{dH_{\text{obs}}}{dt} &= c\gamma\left(H_{\max}-H_{\text{obs}}\right)
  -\dot I_e, \label{eq:H}\\[4pt]
\frac{d\dot I}{dt} &= e\,\frac{dH_{\text{obs}}}{dt}
  +f\mathcal{P}_{\text{input}}-g\dot I. \label{eq:I}
\end{align}

with the v1.0 dimensional assignments
\begin{equation}
[a]=[f]=\mathrm{J^{-1}s^{-1}},\qquad
[b]=[e]=[g]=\mathrm{s^{-1}},\qquad
[c]=1 .
\end{equation}

We treat $\mathcal{E}_{\text{avail}}$, $\mathcal{P}_{\text{input}}$ and
$\dot I_e$ as constant control parameters, and all coefficients
$a,b,c,e,f,g,\gamma_{\max},H_{\max}$ as positive. It is convenient to
introduce the \emph{informational gap}
\begin{equation}
D \equiv H_{\max}-H_{\text{obs}},
\end{equation}
in terms of which \eqref{eq:H} reads $\dot D = -c\gamma D+\dot I_e$.

\section{Equilibrium structure}

\begin{proposition}[Unique equilibrium in closed form]
\label{prop:fp}
For $\dot I_e>0$ the system \eqref{eq:gamma}--\eqref{eq:I} admits exactly one
equilibrium, given by
\begin{equation}
\boxed{\;
\dot I^{*}=\frac{f\mathcal{P}_{\text{input}}}{g},\qquad
\gamma^{*}=\frac{K\gamma_{\max}}{b\gamma_{\max}+K},\qquad
H^{*}_{\text{obs}}=H_{\max}-\frac{\dot I_e}{c\,\gamma^{*}},\;}
\label{eq:fp}
\end{equation}
where $K\equiv a\,\mathcal{E}_{\text{avail}}\dot I^{*}$. Moreover
$0<\gamma^{*}<\gamma_{\max}$.
\end{proposition}

\begin{proof}
At equilibrium $dH_{\text{obs}}/dt=0$, so the first term of \eqref{eq:I}
vanishes identically and \eqref{eq:I} reduces to
$f\mathcal{P}_{\text{input}}-g\dot I^{*}=0$, giving $\dot I^{*}$ as stated.
Substituting into \eqref{eq:gamma} with $d\gamma/dt=0$,
\begin{equation}
K\left(1-\frac{\gamma^{*}}{\gamma_{\max}}\right)=b\gamma^{*}
\;\Longrightarrow\;
K=\gamma^{*}\!\left(b+\frac{K}{\gamma_{\max}}\right)
\;\Longrightarrow\;
\gamma^{*}=\frac{K\gamma_{\max}}{b\gamma_{\max}+K},
\end{equation}
which is a single root, positive, and strictly below $\gamma_{\max}$ for
finite $K$. Finally $dH_{\text{obs}}/dt=0$ gives
$c\gamma^{*}D^{*}=\dot I_e$, hence $D^{*}=\dot I_e/(c\gamma^{*})$ and
$H^{*}_{\text{obs}}=H_{\max}-D^{*}$.
\end{proof}

\begin{corollary}[The entropy sector does not set the processing rate]
\label{cor:decouple}
$\dot I^{*}$ depends only on $f$, $\mathcal{P}_{\text{input}}$ and $g$. It is
independent of $a,b,c,e,\gamma_{\max},H_{\max}$ and $\dot I_e$.
\end{corollary}

\begin{remark}[Why no ignition threshold exists]
The three equations form a closed cycle: $\dot I$ drives $\gamma$ through
\eqref{eq:gamma}, $\gamma$ drives $H_{\text{obs}}$ through \eqref{eq:H}, and
$H_{\text{obs}}$ feeds back into $\dot I$ through \eqref{eq:I}. One might
therefore expect a loop-gain condition separating a quiescent branch from a
self-sustaining branch, i.e.\ bistability with an ignition threshold. No such
structure is present. The reason is that the return path in \eqref{eq:I}
couples to $dH_{\text{obs}}/dt$ rather than to $H_{\text{obs}}$: a derivative
coupling, which vanishes identically at any steady state. The loop is
therefore open at equilibrium and closed only in the transient regime. This is
a structural feature of the ansatz as written, not an artifact of parameter
choice.
\end{remark}

\section{Admissibility constraints}

Proposition~\ref{prop:fp} yields two constraints not stated in v1.0.

\begin{proposition}[Erasure sets the equilibrium gap]
\label{prop:gap}
At equilibrium,
\begin{equation}
H_{\max}-H^{*}_{\text{obs}}=\frac{\dot I_e}{c\,\gamma^{*}} .
\end{equation}
Consequently the v1.0 boundary condition $H_{\text{obs}}\le H_{\max}$ is
satisfied with strict inequality whenever $\dot I_e>0$, and is saturated
exactly, $H^{*}_{\text{obs}}=H_{\max}$, in the limit $\dot I_e\to 0$.
\end{proposition}

Within the ansatz, irreversible erasure is therefore the sole mechanism holding
the equilibrium state below the saturation parameter, and the equilibrium
distance from saturation is inversely proportional to the relaxation rate
$\gamma^{*}$. Under Remark~\ref{rem:hmax} this is a statement about the model's
internal bookkeeping and carries no horizon interpretation.

\begin{proposition}[Admissibility bound on the erasure rate]
\label{prop:admiss}
Since $H_{\text{obs}}$ is a von Neumann entropy in bits, $H_{\text{obs}}\ge 0$
is required. The equilibrium \eqref{eq:fp} is physically admissible only if
\begin{equation}
\boxed{\;\dot I_e\;\le\;c\,\gamma^{*}H_{\max}\;}
\end{equation}
Erasure rates above this bound drive the equilibrium value of
$H_{\text{obs}}$ negative and place the model outside its own domain of
interpretation.
\end{proposition}

\section{Linear stability}

\subsection{Jacobian and two identities}

Linearising \eqref{eq:gamma}--\eqref{eq:I} about \eqref{eq:fp} in the
coordinates $(\gamma,D,\dot I)$ and writing
\begin{equation}
\alpha\equiv\frac{a\mathcal{E}_{\text{avail}}\dot I^{*}}{\gamma_{\max}}+b,
\qquad
\beta\equiv a\mathcal{E}_{\text{avail}}\!\left(1-\frac{\gamma^{*}}{\gamma_{\max}}\right),
\qquad
\mu\equiv c\gamma^{*},
\qquad
\delta\equiv cD^{*},
\end{equation}
all strictly positive, the Jacobian is
\begin{equation}
J=\begin{pmatrix}
-\alpha & 0 & \beta\\
-\delta & -\mu & 0\\
e\delta & e\mu & -g
\end{pmatrix}.
\label{eq:jac}
\end{equation}

Two exact identities simplify everything that follows:
\begin{equation}
\boxed{\;\alpha\beta=a\,\mathcal{E}_{\text{avail}}\,b,
\qquad
\mu\delta=c\,\dot I_e\;}
\label{eq:ident}
\end{equation}
The first follows from
$\alpha=(K+b\gamma_{\max})/\gamma_{\max}$ and
$\beta=a\mathcal{E}_{\text{avail}}b\gamma_{\max}/(b\gamma_{\max}+K)$; the
second from $\delta=cD^{*}=\dot I_e/\gamma^{*}$ and $\mu=c\gamma^{*}$.

Note that $\alpha$, $\beta$ and $\mu$ depend on $\dot I_e$ not at all: by
Proposition~\ref{prop:fp}, neither $\dot I^{*}$ nor $\gamma^{*}$ involves
$\dot I_e$. The erasure rate enters the linearisation through $\delta$ alone,
and does so linearly. This observation drives Proposition~\ref{prop:scaling}
below.

\subsection{Characteristic polynomial}

Expanding $\det(J-\lambda I)$ from \eqref{eq:jac}, the off-diagonal
structure causes the $\mu$-dependent contributions of the third row to cancel
against one another, leaving
\begin{equation}
(\alpha+\lambda)(\mu+\lambda)(g+\lambda)=\beta e\delta\,\lambda,
\end{equation}
that is
\begin{equation}
\lambda^{3}
+\underbrace{(\alpha+\mu+g)}_{a_2}\lambda^{2}
+\underbrace{(\alpha\mu+\alpha g+\mu g-e\beta\delta)}_{a_1}\lambda
+\underbrace{\alpha\mu g}_{a_0}=0 .
\label{eq:cubic}
\end{equation}

\subsection{The Hopf threshold}

\begin{proposition}[Single stability boundary]
\label{prop:hopf}
The equilibrium \eqref{eq:fp} is linearly asymptotically stable if and only if
$e<e_c$, where
\begin{equation}
\boxed{\;
e_c=\frac{(\alpha+\mu+g)(\alpha\mu+\alpha g+\mu g)-\alpha\mu g}
{\beta\delta\,(\alpha+\mu+g)}\;}
\label{eq:ecrit}
\end{equation}
At $e=e_c$ a complex-conjugate pair crosses the imaginary axis at frequency
\begin{equation}
\omega_c=\sqrt{\frac{\alpha\mu g}{\alpha+\mu+g}} .
\label{eq:omega}
\end{equation}
\end{proposition}

\begin{proof}
For the cubic \eqref{eq:cubic} the Routh--Hurwitz criterion requires
$a_2>0$, $a_0>0$ and $a_2a_1>a_0$. The first two hold identically since
$\alpha,\mu,g>0$. The third,
$(\alpha+\mu+g)(\alpha\mu+\alpha g+\mu g-e\beta\delta)>\alpha\mu g$, is linear
and strictly decreasing in $e$ because $\beta,\delta>0$; solving for equality
gives \eqref{eq:ecrit}. At equality $a_1=a_0/a_2$, so \eqref{eq:cubic}
factorises as $(\lambda+a_2)(\lambda^{2}+a_0/a_2)=0$, giving a real root
$-a_2$ and a pair $\pm i\sqrt{a_0/a_2}$, which is \eqref{eq:omega}. Since only
$a_1$ depends on $e$ and does so monotonically, the crossing is unique and
transversal.
\end{proof}

\begin{corollary}[Zero-erasure limit is unconditionally stable]
\label{prop:scaling}
Because $\alpha,\beta,\mu$ are independent of $\dot I_e$ and
$\delta=\dot I_e/\gamma^{*}$, the product $e_c\dot I_e$ is invariant under
changes of $\dot I_e$:
\begin{equation}
e_c\;\propto\;\frac{1}{\dot I_e},
\qquad\text{equivalently}\qquad
e_c=\frac{\alpha\mu\left[(\alpha+\mu+g)(\alpha\mu+\alpha g+\mu g)-\alpha\mu g\right]}
{a\,\mathcal{E}_{\text{avail}}\,b\,c\,\dot I_e\,(\alpha+\mu+g)} ,
\end{equation}
using \eqref{eq:ident}. Hence $e_c\to\infty$ as $\dot I_e\to 0$: in the
zero-erasure limit the equilibrium is linearly stable for every finite
feedback gain $e$, and the Hopf frequency \eqref{eq:omega} is unchanged.
\end{corollary}

\begin{remark}[Interpretation, stated narrowly]
Within this ansatz the only linear instability available to the system is
driven by the irreversible-erasure channel. This is of interest because v1.0
explicitly declines to specify any relation between $\dot I_{\text{model}}$ and
$\dot I_{\text{erase}}$ and records that omission as an open problem. The
present result sharpens what is at stake in that omission: the unspecified
relation controls not merely a dissipative cost but the existence of
oscillatory dynamics. It does not supply the missing relation, and no relation
should be inferred from it.
\end{remark}

\section{Numerical verification}

All results were checked against direct numerical evaluation using the
benchmark parameter set
$a=b=c=f=g=\mathcal{E}_{\text{avail}}=1$,
$\gamma_{\max}=10$, $H_{\max}=100$, $\mathcal{P}_{\text{input}}=5$,
in the units of v1.0.

\begin{table}[h]
\centering
\begin{tabularx}{\textwidth}{Xccc}
\toprule
\textbf{Quantity} & $\dot I_e=1$ & $\dot I_e=10$ & $\dot I_e=37$\\
\midrule
$\gamma^{*}$ (s$^{-1}$), Eq.~\eqref{eq:fp} & $3.3333$ & $3.3333$ & $3.3333$\\
$\dot I^{*}$ (s$^{-1}$), Eq.~\eqref{eq:fp} & $5.0000$ & $5.0000$ & $5.0000$\\
$H^{*}_{\text{obs}}$ (bits), Eq.~\eqref{eq:fp} & $99.7000$ & $97.0000$ & $88.9000$\\
Residual $\max|\dot{\mathbf{x}}(\mathbf{x}^{*})|$ & $9.8\times10^{-15}$ & $4.4\times10^{-16}$ & ---\\
$e_c$ analytic, Eq.~\eqref{eq:ecrit} & $44.880952$ & $4.488095$ & $1.212999$\\
$e_c$ numeric (eigenvalue crossing) & $44.880952$ & $4.488095$ & $1.212999$\\
$e_c\,\dot I_e$ (invariant, Cor.~\ref{prop:scaling}) & $44.880952$ & $44.880952$ & $44.880952$\\
$\omega_c$ analytic, Eq.~\eqref{eq:omega} & $0.925820$ & $0.925820$ & $0.925820$\\
$\omega_c$ numeric & $0.925820$ & $0.925820$ & $0.925820$\\
$\max\operatorname{Re}\lambda$ at $e=e_c$ & $-1.5\times10^{-16}$ & $9.8\times10^{-16}$ & $-3.4\times10^{-16}$\\
\bottomrule
\end{tabularx}
\caption{Analytic predictions against numerical evaluation. Agreement is to
machine precision throughout.}
\end{table}

\section{Global behaviour and a required domain restriction}

Linear stability is local. Direct integration of
\eqref{eq:gamma}--\eqref{eq:I} shows that the basin of attraction is bounded
and that trajectories leaving it do not merely diverge but reach
\emph{unphysical} states in finite time.

The mechanism is explicit in the equations. If $\gamma$ is driven negative
while the gap $D$ remains positive, \eqref{eq:H} makes $dH_{\text{obs}}/dt$
strongly negative; \eqref{eq:I} then drives $\dot I$ negative; and with
$\dot I<0$ and $(1-\gamma/\gamma_{\max})>1$ the right-hand side of
\eqref{eq:gamma} is negative and growing in magnitude, so $\gamma$ runs away.
Nothing in the ansatz prevents this: the system has no positivity-preserving
structure for $\gamma$ or $\dot I$.

Numerically, with $\dot I_e=10$ and an initial perturbation of $10^{-4}$ in
$\gamma$:

\begin{table}[h]
\centering
\begin{tabularx}{\textwidth}{lX}
\toprule
\textbf{Gain} & \textbf{Late-time behaviour}\\
\midrule
$e=0.90\,e_c$ & returns to \eqref{eq:fp}; oscillation amplitude $<10^{-7}$\\
$e=1.01\,e_c$ & bounded oscillation, $\gamma$ amplitude $5.4\times10^{-3}$\,s$^{-1}$\\
$e=1.02\,e_c$ & bounded oscillation, $\gamma$ amplitude $4.9\times10^{-2}$\,s$^{-1}$\\
$e=1.04\,e_c$ & escapes: $\gamma\to 0^{-}$ at $t\approx259$\\
$e=1.08\,e_c$ & escapes: $\gamma\to 0^{-}$ at $t\approx136$\\
$e=1.20\,e_c$ & escapes: $\gamma\to 0^{-}$ at $t\approx55$\\
\bottomrule
\end{tabularx}
\caption{Behaviour across the stability boundary. A bounded oscillatory regime
exists in a narrow window above $e_c$; beyond it the trajectory leaves the
physical domain in finite time.}
\end{table}

Accordingly the ansatz requires an explicit statement of its domain of
validity. Two minimal options, either of which is a modelling choice rather
than a derived result:
\begin{enumerate}
\item restrict the admissible parameter range to $e<e_c$ with
Proposition~\ref{prop:admiss} enforced; or
\item modify the ansatz so that $\gamma\ge0$ and $\dot I\ge0$ are invariant
(for example by making the decay terms multiplicative).
\end{enumerate}
Option 2 would change the frozen equations and therefore belongs to a future
versioned artifact, not to this appendix.

\section{What this appendix does and does not establish}

\noindent\textbf{Established.} Within the frozen ansatz, and as consequences of
it alone: uniqueness and closed form of the equilibrium
(Prop.~\ref{prop:fp}); independence of the steady-state processing rate from
the entropy sector (Cor.~\ref{cor:decouple}); the erasure-set equilibrium gap
(Prop.~\ref{prop:gap}); the admissibility bound (Prop.~\ref{prop:admiss}); a
unique Hopf boundary with closed-form threshold and frequency
(Prop.~\ref{prop:hopf}); the exact inverse scaling of that threshold with the
erasure rate (Cor.~\ref{prop:scaling}); and the finite-time escape requiring a
domain restriction.

\medskip
\noindent\textbf{Not established.} This appendix does not derive the ansatz
from microscopic non-equilibrium dynamics; it analyses a system that v1.0
already labels phenomenological. It does not relate $\dot I_{\text{model}}$ to
$\dot I_{\text{erase}}$. It does not give $H_{\max}$ any geometric or horizon
meaning (Remark~\ref{rem:hmax}). It does not bear on the Gate~0 question, and
its results must not be cited as progress on the observer--horizon mapping,
which remains \textbf{UNRESOLVED IN GENERALITY}. It makes no cosmological
prediction and does not constrain $(\epsilon,\beta,\phi_0)$ in v1.0's
gravitational-wave template. It determines no numerical values for
$a,\dots,g$; the benchmark set above is illustrative only, chosen to exhibit
the structure, and the qualitative conclusions are parameter-independent while
the tabulated numbers are not.

\medskip
\noindent\textbf{Nearest open problems made concrete by this analysis.}
(i) Specify the $\dot I_{\text{model}}$--$\dot I_{\text{erase}}$ relation and
recompute $e_c$, which by Cor.~\ref{prop:scaling} then ceases to be a free
parameter. (ii) Determine the criticality of the Hopf bifurcation
analytically; the numerics above are consistent with a supercritical
bifurcation whose limit cycle is rapidly destroyed, but no normal-form
computation is offered here. (iii) Establish whether any positivity-preserving
modification reproduces the same $e_c$ scaling.

\section{Reproducibility}

All numerical results were produced by
\texttt{physics\_v1\_0\_ode\_analysis.py}, accompanying this appendix, using
NumPy 2.4.4 and SciPy 1.17.1. Equilibria are evaluated in closed form and
verified by residual; stability boundaries are located both analytically via
\eqref{eq:ecrit} and numerically by bracketing the sign change of
$\max\operatorname{Re}\lambda(J)$; trajectories are integrated with LSODA at
tolerances $\mathrm{rtol}=10^{-9}$, $\mathrm{atol}=10^{-12}$, with a
terminal event on $\gamma=0$. File hashes are recorded in the accompanying
manifest.

\end{document}
